% !PRINTED 119
\weberpage{119}
By repeated application of (9) it follows, if $a$, $b$ are integers [cf. \S\,17, (16)]:

\weq{11}{\sigma(u+a\omega_1+b\omega_2) = (-1)^{a+b+ab}\,e^{(a\eta_1+b\eta_2)(2u+a\omega_1+b\omega_2)}\,\sigma(u).}
If one develops the function $\sigma(u)$, as it is given by formula (6), in powers of $u$, it is found that not only the second but also the third power of $u$ fails to occur in the development, and one has therefore:

\weq{12}{\sigma(0)=0,\quad \sigma'(0)=1,\quad \sigma''(0)=0,\quad \sigma'''(0)=0.}
If to $\omega_1, \omega_2$ a linear substitution

\weq{13}{(\omega_1', \omega_2') = \begin{pmatrix}\alpha, & \beta\\ \gamma, & \delta\end{pmatrix}(\omega_1, \omega_2)}
is applied, then the quantities $\eta_1, \eta_2$ undergo the corresponding substitution

\weq{14}{(\eta_1', \eta_2') = \begin{pmatrix}\alpha, & \beta\\ \gamma, & \delta\end{pmatrix}(\eta_1, \eta_2),}
as one infers immediately from the relations (7) and (11).

If the formulae (9) are differentiated logarithmically with respect to $u$, there follows

\weq{15}{\begin{aligned}
\frac{\sigma'(u+\omega_1)}{\sigma(u+\omega_1)} &= \frac{\sigma'(u)}{\sigma(u)} + 2\eta_1,\\
\frac{\sigma'(u+\omega_2)}{\sigma(u+\omega_2)} &= \frac{\sigma'(u)}{\sigma(u)} + 2\eta_2,
\end{aligned}}
and by setting, in the first, $u=-\dfrac{\omega_1}{2}$, in the second $u=-\dfrac{\omega_2}{2}$, and observing that $\sigma(u)$ is an odd and consequently $\sigma'(u)$ an even function of $u$:

\weq{16}{\eta_1 = \frac{\sigma'\left(\dfrac{\omega_1}{2}\right)}{\sigma\left(\dfrac{\omega_1}{2}\right)}.\quad \eta_2 = \frac{\sigma'\left(\dfrac{\omega_2}{2}\right)}{\sigma\left(\dfrac{\omega_2}{2}\right)}.}

\begin{center}
\textbf{\S\,36. \gesperrt{The Functions $\sigma_{00}, \sigma_{01}, \sigma_{10}$.}}
\end{center}

It still remains for us to obtain the analogous results for the remaining characteristics. We start from the following remark: if $x_1, x_2$; $y_1, y_2$ are any two pairs of quantities which are transformed by one and the same linear substitution $S$ into $x_1', x_2'$; $y_1', y_2'$, then we have:
%CONT>>
